%% kapa: exam=1 round=22 year=2011 period=1 form=A subject=economics subject_ko=경제원론 seq=29 number=69
%% source: Q-Net 감정평가사 기출문제 (cst00309) artlSeq=5036790 file=extracted/1008/5036790/2011년 제22회 감정평가사 1차시험 기출문제/1교시.hwp
%% conversion: kordoc 4.24.51 -> tools/gen1.py (auto); review=proofread-sonnet; flags=none
\begin{question}{69}
전통적 화폐수량설에 근거한 화폐의 중립성이 성립할 경우 다음 설명 중 옳지 않은 것은?
\begin{choices}
  \item 통화량 증가율을 증가시키면 명목이자율이 상승한다.
  \item 통화량 증가율을 증가시키면 인플레이션율이 상승한다.
  \item 통화량을 증가시켜도 실질 국민소득수준은 변화하지 않는다.
  \item 통화량을 증가시키면 실업률은 하락한다.
  \item 통화량을 증가시켜도 실질이자율은 변화하지 않는다.
\end{choices}
\end{question}
\answer{4}
